\subsection{A weaker criterion from the deck action on two-primary torsion}

\begin{claim}{A.C2.7}{Proved}
Let $C$ be a smooth proper geometrically connected curve of genus at
least two over a field $K$ finitely generated over $\F_2$. Suppose that,
over $\overline K$, it admits a connected nontrivial finite \'etale
Galois cover $D\to C$ of odd degree, with group $G$, for which
\[
             G\longrightarrow
             \operatorname{Aut}_{\F_2}(J_D[2](\overline K))
                                                               \tag{C2.19}
\]
is not faithful. Then the geometrically prime-to-two point-section
map of $C$ is injective. The conclusion also holds for its
finite-extension forms and hyperbolic open subcurves.

More precisely, a pair of distinct points giving equal section classes
would force (C2.19) to be faithful for every such cover. If
$m_2(G)$ denotes the least dimension of a faithful representation of
$G$ over $\F_2$, the same hypothetical pair would force
\[
                            f(D)-f(C)\ge m_2(G).               \tag{C2.20}
\]
In particular, a single such cover with $f(D)=f(C)$ suffices for
point-section injectivity.
\end{claim}
\begin{proof}
Work first over an algebraic closure. The curve $D$ is not
hyperelliptic. Otherwise its hyperelliptic involution is unique,
so every deck transformation descends to its genus-zero quotient.
Choose an element $\sigma\in G$ of odd prime order. Its image on
$\mathbf P^1$ has the same order, since the kernel of descent is
the order-two hyperelliptic group. A finite odd-order automorphism
of $\mathbf P^1$ in characteristic two has a fixed geometric
point. Its fibre on $D$ has at most two geometric points, so
$\sigma$ fixes a point in that fibre. This contradicts the free
deck action of an \'etale Galois cover.

Let $P',Q'\in D(\overline K)$ lie over distinct points of $C$.
The class $T=[Q'-P']$ has trivial stabilizer in $G$. Indeed,
$\delta T=T$ for $1\ne\delta\in G$ would give
\[
                     \delta Q'+P'\sim Q'+\delta P'.
                                                               \tag{C2.21}
\]
The two effective degree-two divisors in (C2.21) are distinct:
their terms lying over each of the two base points would
otherwise force both $\delta P'=P'$ and $\delta Q'=Q'$.
A linear equivalence between distinct effective degree-two
divisors on a curve of genus at least two gives a degree-two
map to $\mathbf P^1$, after removing any common base point.
A degree-one map is excluded by the genus. Thus (C2.21) would
make $D$ hyperelliptic, a contradiction.

The \'etale two-adic Tate module of $J_D$ is free of rank $f(D)$,
and its geometric two-primary torsion is the quotient of its
rational Tate module by this lattice. An odd-order automorphism
acting trivially on $J_D[2](\overline K)$ acts trivially on the
whole two-adic module. This can be checked without any
torsion-free congruence-group citation: if a matrix $A$ of odd
order $m$ satisfies $A\equiv1\pmod{2^r}$, then
\[
 A^m\equiv1+m(A-1)\pmod{2^{r+1}}
\]
implies $A\equiv1\pmod{2^{r+1}}$. Starting at $r=1$ and
iterating proves $A=1$. Thus any element in the kernel of
(C2.19) fixes every two-primary torsion point of $J_D$.

Now suppose that distinct $P,Q\in C(K)$ give equal section
classes. Descend the cover and its deck transformations after
a finite extension of $K$, and extend again to make a chosen
lift $P'$ rational. By A.3 there is a lift $Q'$ with
$T=[Q'-P']$ two-primary. A nonidentity element in the kernel
of (C2.19) would fix $T$, contrary to its trivial stabilizer.
This proves the criterion and the necessity of faithfulness.
The same argument after restriction proves the stated form and
open-subcurve variants.

For the quantitative assertion, the averaging idempotent
$|G|^{-1}\sum_{g\in G}g$ acts integrally on the two-adic Tate
module, since $|G|$ is odd. Its invariant summand is the
pullback of the two-adic Tate module of $J_C$: pullback
followed by norm is multiplication by $|G|$, and norm followed
by pullback is the sum of deck transformations. This summand
has rank $f(C)$. The complementary summand has rank
$f(D)-f(C)$, and its reduction modulo two must be a faithful
$G$-representation because the invariant summand is trivial.
This proves (C2.20).
\end{proof}

\begin{claim}{A.C2.8}{Proved}
In the genus-two cyclic situation of A.C2.5, let $S$ be the
union of the nonzero doubling orbits that miss $\Theta$.
The kernel of the deck action on geometric two-primary torsion
has order
\[
                    d=\gcd\bigl(n,\{i:i\in S\}\bigr),        \tag{C2.22}
\]
where $d=n$ if $S$ is empty. Consequently $d>1$ already implies
point-section injectivity by A.C2.7.

In particular, if $n=\ell^a$ is an odd prime power, it suffices
that every doubling orbit on $(\mathbf Z/n)^\times$ meets
$\{i:L^i\in\Theta\}$. No condition is needed on characters
of smaller order. The simple sufficient condition is
\[
 L\in\Theta\text{ of exact order }\ell^a,\qquad
                 \langle2,-1\rangle=(\mathbf Z/\ell^a)^\times.
                                                               \tag{C2.23}
\]
For a cyclic group of odd prime-power order $\ell^a$ one has
$m_2(G)=\operatorname{ord}_{\ell^a}(2)$.
\end{claim}
\begin{proof}
The Artin--Schreier description of \'etale $H^1$ identifies its
scalar extension with the stable Frobenius part of
$H^1(D,\mathcal O_D)$. The calculation of A.C2.5 therefore
shows that its nontrivial cyclic characters are exactly those
indexed by $S$, with multiplicity one. Dualizing to geometric
two-torsion may invert all characters; it does not change
their common kernel. A power $\delta^k$ acts trivially exactly
when $n\mid ik$ for every $i\in S$. The number of such powers
is the greatest common divisor in (C2.22).

If $n=\ell^a$, this kernel is nontrivial exactly when no
primitive character survives, since a primitive index is
coprime to $n$, whereas all nonprimitive indices are divisible
by $\ell$. Symmetry of $\Theta$ gives (C2.23) as in A.C2.5.
For example at $n=9$, a theta point of exact order nine kills
the single orbit of all six primitive characters. The
order-three block may remain ordinary, but the deck action
then has a kernel of order three, which suffices here. This
does not assert existence of such a nine-torsion theta point
on every curve.

Finally, an irreducible $\F_2$-representation of a cyclic
group corresponds to a doubling orbit of its geometric
characters. A faithful character of order $\ell^a$ has orbit
length $\operatorname{ord}_{\ell^a}(2)$. Any faithful
representation of this cyclic prime-power group must contain
such a character: otherwise all character orders divide
$\ell^{a-1}$ and a subgroup of order $\ell$ acts trivially.
One faithful orbit already gives a faithful representation,
proving the formula for $m_2(G)$.
\end{proof}

The full-quotient obstruction in A.C2.6 and the weaker deck
criterion are compatible. Surviving order-three characters
prevent the entire Prym from having $2$-rank zero, but at
degree $3^a$, $a>1$, they need not make the deck action
faithful. Torsion density still does not guarantee the
prime-power order or simultaneous primitive-orbit
intersections required by (C2.23). A degree-sized lower
bound in place of (C2.20), or a construction forcing a
sufficiently small normalized $2$-rank in general towers,
has not been established here.
